\begin{figure}[H]
  \centering

  \begin{tikzpicture}[
    scale=0.8,
    every node/.style={scale=0.8},
    a/.style={
      font=\bfseries,
      text=DarkGreen,
    },
    b/.style={
      % font=\bfseries,
      % text=DarkRed,
    },
    bit/.style={
      draw=none,
      font=\scriptsize,
      text=gray,
    },
    ]

    \path[
    % Style for tree only.
    level distance=20mm,
    level 1/.style={sibling distance=70mm},
    level 2/.style={sibling distance=50mm},
    level 3/.style={sibling distance=30mm},
    every node/.append style={
      circle,
      draw,
      inner sep=2pt,
      minimum size=10pt,
    }]
    node (root) {}
    child {node (n0) {}
      child {node (n00) {}
        child {node (n000) {}
          child {node {}
            child {node (n00011) {}
              edge from parent node[bit, left] {1}
            }
            edge from parent node[bit, left] {1}
          }
          edge from parent node[bit, above left] {0}
        }
        child {node {}
          child {node {}
            child {node (n00101) {}
              edge from parent node[bit, left] {1}
            }
            edge from parent node[bit, above left] {0}
          }
          child {node {}
            child {node (n00110) {}
              edge from parent node[bit, left] {0}
            }
            edge from parent node[bit, below left] {1}
          }
          edge from parent node[bit, below left] {1}
        }
        edge from parent node[bit, above left] {0}
      }
      child {node (n01) {}
        child {node (n010) {}
          child {node {}
            child {node (n01001) {}
              edge from parent node[bit, left] {1}
            }
            edge from parent node[bit, left] {0}
          }
          edge from parent node[bit, left] {0}
        }
        edge from parent node[bit, below left] {1}
      }
      edge from parent node[bit, above left] {0}
    }
    child {node (n1) {}
      child {node (n10) {}
        child {node {}
          child {node {}
            child {node (n10100) {}
              edge from parent node[bit, left] {0}
            }
            edge from parent node[bit, left] {0}
          }
          edge from parent node[bit, left] {1}
        }
        edge from parent node[bit, left] {0}
      }
      edge from parent node[bit, below left] {1}
    };

    \node[a, below=5pt of n00011] {3};
    \node[a, below=16pt of n00011] {$00011_{(2)}$};
    \node[a, below=5pt of n00101] {5};
    \node[a, below=16pt of n00101] {$00101_{(2)}$};
    \node[a, below=5pt of n00110] {6};
    \node[a, below=16pt of n00110] {$00110_{(2)}$};
    \node[a, below=5pt of n01001] {9};
    \node[a, below=16pt of n01001] {$01001_{(2)}$};
    \node[a, below=5pt of n10100] {20};
    \node[a, below=16pt of n10100] {$10100_{(2)}$};

    \node[b, left=10pt of root] {
      $3=00011_{(2)}$,
      $7=00111_{(2)}$,
      $12=01100_{(2)}$,
      $13=01101_{(2)}$,
      $19=10011_{(2)}$
    };
    \node[b, left=10pt of n0] {
      $00011_{(2)}$,
      $00111_{(2)}$,
      $01100_{(2)}$,
      $01101_{(2)}$
    };
    \node[b, left=10pt of n1] {
      $10011_{(2)}$
    };
    \node[b, left=2pt of n00] {
      $01100_{(2)}$,
      $01101_{(2)}$
    };
    \node[b, left=2pt of n01] {
      $00011_{(2)}$,
      $00111_{(2)}$
    };
    \node[b, left=2pt of n000] {
      $01100_{(2)}$,
      $01101_{(2)}$
    };
    \node[b, left=2pt of n010] {
      $00011_{(2)}$
    };
  \end{tikzpicture}

  \caption{Trie-ul construit pentru exemplul 2 din enunț, în care dorim să obținem al $k=10$-lea cel mai mic xor. În rădăcină obținem $4 \cdot 4 + 1 \cdot 1 = 17$ perechi cu xor-ul 0, deci primul bit al răspunsului este 0. Pe nivelul 1 (corespunzător celui de-al doilea bit din stînga) obținem $2 \cdot 3 + 2 \cdot 1 + 1 \cdot 1 = 9$ perechi cu xor-ul 0, deci al doilea bit al rezultatului este 1.}
  \label{fig:trie-esoteric-bovines}
\end{figure}
